\documentclass[11pt]{article}
\usepackage[margin=1in]{geometry}
\usepackage{amsmath,amssymb}
\usepackage[hidelinks]{hyperref}

\title{Disproof of Conjecture 00000004103\\
\large A symmetric star with a torsion-free graph braid group}
\author{}
\date{}

\begin{document}
\maketitle

Let \(S=K_{1,4}\) be the star with center \(c\) and leaves
\(0,1,2,3\). The map fixing \(c\) and sending each leaf \(i\) to
\(3-i\) is an involutive graph automorphism. It exchanges two pairs
of arms and fixes only the center. Nevertheless its two-strand
graph braid group has no element of order two.

\section*{The two-strand braid group of \(S\)}

Let \(U_2(S)\) be the space of unordered pairs of distinct points of
the geometric graph \(S\), so that \(B_2(S)=\pi_1(U_2(S))\). The
unordered discrete configuration space \(UD_2(S)\) consists of pairs
of cells of \(S\) with disjoint closures, with the two coordinates
identified under swapping.

The stable equivalence theorem of Prue and Scrimshaw
\cite[Theorem~3.2]{prue-scrimshaw}, including its unordered version,
applies here. For \(n=2\), the graph must be finite and connected
with at least two vertices. The required distances between distinct
essential vertices are at least \(n-1=1\) edge, and homotopically
essential closed paths must have at least \(n+1=3\) edges. In \(S\)
the distance condition holds for any two distinct vertices, and
there are no homotopically essential closed paths because \(S\) is
a tree. Thus \(U_2(S)\) deformation retracts onto \(UD_2(S)\).

The four vertices of \(UD_2(S)\) of the form \(\{c,i\}\) correspond
to the vertices of \(K_4\). The six vertices of the form \(\{i,j\}\),
with \(i\ne j\) both leaves, correspond to the midpoints of its edges.
The one-cell given by a point moving on the arm \(ci\), with a point
fixed at the different leaf \(j\), joins \(\{c,j\}\) to \(\{i,j\}\).
Hence each midpoint has precisely its two expected incident edges.
There are no two-cells because any two closed arms meet at \(c\).
This identifies \(UD_2(S)\) with the barycentric subdivision of \(K_4\).

The resulting graph is connected and has ten vertices and twelve
edges. Contracting a spanning tree gives a wedge of
\(12-10+1=3\) circles, so \(B_2(S)\) is a free group of rank three.
In a free group, any nonidentity element has a nonempty cyclically
reduced conjugate, whose positive powers remain reduced and nonempty.
Thus \(B_2(S)\) is torsion-free. In particular it contradicts the
claim that an involutive star symmetry implies an element of order
two in the graph braid group.

\section*{Lean verification}

The Lean~4.19.0 project checks the star automorphism, the incidences,
the connectivity and cell counts of the finite configuration graph,
and its cycle rank. It also proves directly in the free-group normal
form that a self-inverse reduced word is empty, which excludes order
two. The deformation retraction and the passage from the finite
graph to its free fundamental group are given above. Run
\verb|lake build| in \verb|lean/| to verify the formal results.

\begin{thebibliography}{1}
\bibitem{prue-scrimshaw}
P.~Prue and T.~Scrimshaw,
\emph{Abrams's stable equivalence for graph braid groups},
\href{https://arxiv.org/abs/0909.5511}{arXiv:0909.5511}.
\end{thebibliography}

\end{document}
